\caption{Finite-data recovery of reduced non-normal geometry using the corrected $R_A$ diagnostic. The stable $N$-dimensional VAR(1) generating model has $\rho(F)=0.95$ and white Gaussian noise of standard deviation $0.05$. For each dimension the coupling is recalibrated to obtain $R_{A,\mathrm{ref}}\simeq0.2,1,10$; hence these are not identical generating operators to those in the originally submitted figure. Solid lines denote M2 and dashed lines M3. Blue, orange and green denote the three $R_{A,\mathrm{ref}}$ regimes. Panels (a,b) vary the number of trajectories $M$ ($N=100$, $T_{\mathrm{train}}=200$); (c,d) vary dimension $N$ ($M=25$, $T_{\mathrm{train}}=200$); and (e,f) vary training horizon $T_{\mathrm{train}}$ ($M=25$, $N=100$). The left column shows the largest principal angle to the M2 reference plane obtained from the known operator and, on the grey right axis, the full-operator relative Frobenius error for the $R_{A,\mathrm{ref}}\simeq1$ regime. The right column shows $\mathcal E_{R_A}=|\log_{10}(\widehat R_A/R_{A,\mathrm{ref}})|$, with the horizontal line at $\log_{10}2$. The insets in (c,e) retain the single-trajectory ($M=1$) plane-angle comparisons. The $M=1$ amplification-ratio results that formerly occupied the small insets in (d,f) are instead shown in a separate companion figure because so few are admissible. Undefined $R_A$ estimates are not assigned values or connected across missing points. Every setting uses one archived synthetic noise realization; no cross-realization uncertainty interval is implied. The optional diagnostic companion overlays faint $|\log_{10}(\widehat R_\sigma/R_{\sigma,\mathrm{ref}})|$ curves in the right column.}
